%% toc: P15 159 159.1,159.2
% Independence of families of arbitrary size, and the complement-choice form of mutual independence.
\chapter{Independent Families of Arbitrary Size and the Complement Test}
\label{sec:p15-w02b-p15-006-arbitrary-independence}

Chapter~\ref{sec:p15-independent-random-variables} defines independence
for families indexed by a finite or countably infinite set. Durrett's
author Version~5, dated 11 January 2019, imposes no such restriction:
after defining independence for finitely many sigma-fields, random
variables or sets, it calls an infinite collection of such objects
independent when every finite subcollection is independent
\cite[p.~44]{DurrettPTEAuthor2019}. This chapter states that definition
for an index set of any cardinality, proves that every structural result
of the countable treatment carries over to it, and records what the
uncountable case can and cannot contain on the probability spaces
constructed in this part. It then proves the complement-choice form of
mutual independence for finitely many events, which Grinstead and Snell
give as an equivalent form in their Definition~4.2
\cite[sec.~4.1, p.~141]{GrinsteadSnellChance2006}, together with the
converse that the earlier chapters do not supply.

\emph{Inputs.} A probability space \((\Omega,\mathcal F,\mathbb P)\);
an index set \(I\) of any cardinality; events, collections of events,
sub-sigma-algebras or random elements indexed by \(I\). The finite and
countable results of Chapters~\ref{sec:p15-finite-probability}
and~\ref{sec:p15-independent-random-variables} are used as proved there.
\emph{Conclusions.} The definition of independence for arbitrary
families and its agreement with the countable definitions; transfer of
the generated-information, complement, indicator, null-replacement,
pi-system, threshold, bounded-function and disjoint-group results to
arbitrary families; the probability of a countable intersection of
independent events; the equivalence of the subset form and the
complement-choice form of mutual independence; and a countability
theorem for nondegenerate independent families on any probability space
with a countable approximating collection of events, including the
uniform probability space on \((0,1)\).

\emph{Relationships.} The countable definitions
\cref{def:p15-general-independence,def:p15-independent-random-elements}
are the special cases of
\cref{def:p15-w02b-p15-006-arbitrary-families} with countable \(I\).
The finite event definition \cref{def:p15-finite-independence} is the
case of finitely many events on a finite outcome space. The forward
half of the complement form is the complement clause of
\cref{thm:p15-independence-basic}. The pairwise-but-not-mutual
boundaries of \cref{def:p15-pairwise-not-mutual} and
\cref{thm:p15-independence-counterexamples} are re-examined with the
complement test in \cref{def:p15-w02b-p15-006-complement-examples}.
Infinite product measures over a countable index set are constructed in
Chapter~\ref{sec:p15-countable-product-measures}; no product measure over
an uncountable index set is constructed in this part
(\cref{def:p15-countable-product-next}).

\section{Independence of events, random variables and sigma-algebras over an arbitrary index set}
\label{sec:p15-w02b-p15-006-families}

\subsection{The definition and its finite character}
\label{sec:p15-w02b-p15-006-definition}

\begin{hbdefinition}{def:p15-w02b-p15-006-arbitrary-families}{Independence of a family of any size}
Fix a probability space \((\Omega,\mathcal F,\mathbb P)\) and an
arbitrary set \(I\). It may be empty, finite, countably infinite or
uncountable.
\begin{enumerate}
\item Events \((A_i)_{i\in I}\) in \(\mathcal F\) are
      \emph{independent} if for every finite \(J\subseteq I\),
      \[
       \mathbb P\left(\bigcap_{i\in J}A_i\right)
                              =\prod_{i\in J}\mathbb P(A_i),
      \]
      with \(\bigcap_{i\in\varnothing}A_i=\Omega\) and the empty
      product equal to one.
\item Collections \(\mathcal C_i\subseteq\mathcal F\), \(i\in I\), are
      \emph{independent} if the same factorization holds for every
      finite \(J\subseteq I\) and every selection
      \(A_i\in\mathcal C_i\), \(i\in J\). Sub-sigma-algebras
      \(\mathcal F_i\subseteq\mathcal F\) are independent when they are
      independent as collections.
\item For measurable spaces \((S_i,\mathcal S_i)\), random elements
      \(X_i:\Omega\to S_i\), \(i\in I\), are \emph{independent} if for
      every finite \(J\subseteq I\) and every selection
      \(B_i\in\mathcal S_i\), \(i\in J\),
      \[
       \mathbb P\left(\bigcap_{i\in J}\{X_i\in B_i\}\right)
                       =\prod_{i\in J}\mathbb P(X_i\in B_i).
      \]
      Real random variables are the case
      \((S_i,\mathcal S_i)=(\mathbb R,\mathcal B(\mathbb R))\),
      which is the source's setting.
\end{enumerate}
A \emph{subfamily} is the family indexed by a subset \(J\subseteq I\);
it is \emph{finite} or \emph{countable} with \(J\). A family is
\emph{pairwise independent} when every subfamily indexed by a
two-element set is independent.

Only finite subfamilies enter the definition. No equation is imposed
on an intersection over an infinite subfamily. A sigma-algebra is
closed under countable intersections only, so for an uncountable
subfamily such an intersection is not even required to be an event.
\index{independence!arbitrary index set}
\end{hbdefinition}

\begin{hbtheorem}{thm:p15-w02b-p15-006-finite-character}{proved=annex:p15-w02b-p15-006-finite-character}{Finite character and countable intersections}
\par\nobreak
Let \(I\) be arbitrary.
\begin{enumerate}
\item For a family of events, of collections of events, or of random
      elements indexed by \(I\), the following are equivalent:
      the family is independent; every finite subfamily is independent;
      every countable subfamily is independent in the sense of
      \cref{def:p15-general-independence} or
      \cref{def:p15-independent-random-elements}. Every subfamily of an
      independent family is independent.
\item When \(I\) is finite or countably infinite,
      \cref{def:p15-w02b-p15-006-arbitrary-families} coincides with
      \cref{def:p15-general-independence,def:p15-independent-random-elements}.
\item Let \((A_i)_{i\in I}\) be independent events and let
      \(J\subseteq I\) be countably infinite, enumerated without
      repetition as \(j_1,j_2,\ldots\). The partial products
      \(\prod_{k=1}^n\mathbb P(A_{j_k})\) are nonincreasing in
      \(n\), and
      \[
       \mathbb P\left(\bigcap_{i\in J}A_i\right)
             =\lim_{n\to\infty}\prod_{k=1}^n\mathbb P(A_{j_k}).
      \]
      In particular the limit does not depend on the enumeration.
\end{enumerate}
\end{hbtheorem}

\begin{hbtheorem}{thm:p15-w02b-p15-006-arbitrary-transfer}{proved=annex:p15-w02b-p15-006-arbitrary-transfer}{The countable-family results for arbitrary families}
\par\nobreak
Let \(I\) be arbitrary.
\begin{enumerate}
\item Random elements \(X_i\), \(i\in I\), are independent exactly
      when the sigma-algebras \(\sigma(X_i)\) are independent. If
      independent sub-sigma-algebras \(\mathcal F_i\) make each
      \(X_i\) measurable, the \(X_i\) are independent. Replacing
      independent sigma-algebras by smaller ones preserves independence.
      For measurable maps \(f_i:(S_i,\mathcal S_i)\to(T_i,\mathcal T_i)\),
      independence of the \(X_i\) implies independence of the
      \(f_i(X_i)\).
\item Let \((A_i)_{i\in I}\) be independent events and \(K\subseteq I\)
      any subset, finite or infinite. Setting \(B_i=A_i^c\) for
      \(i\in K\) and \(B_i=A_i\) otherwise gives independent events
      \((B_i)_{i\in I}\). Events \(A_i\) are independent exactly when
      their indicators are independent, and exactly when the
      sigma-algebras \(\sigma(\{A_i\})=\{\varnothing,A_i,A_i^c,\Omega\}\)
      are independent.
\item If real \(X_i\), \(i\in I\), are independent and each real
      random variable \(Y_i\) equals \(X_i\) outside a measurable null
      set \(N_i\), then the \(Y_i\) are independent and \(Y_i\) has the
      law of \(X_i\). No condition is placed on
      \(\bigcup_{i\in I}N_i\), which for uncountable \(I\) need not be
      null.
\item If collections \(\mathcal C_i\), \(i\in I\), are independent and
      each is a pi-system, then the \(\sigma(\mathcal C_i)\) are
      independent.
\item Real random variables \(X_i\), \(i\in I\), are independent
      exactly when, for every finite list \(i_1,\ldots,i_n\) of
      distinct indices,
      \(\mathbb P(X_{i_1}\leq t_1,\ldots,X_{i_n}\leq t_n)
      =\prod_{k=1}^n\mathbb P(X_{i_k}\leq t_k)\)
      for all \(t_k\in\mathbb Q\cup\{\infty\}\), with
      \(\{X\leq\infty\}=\Omega\).
\item Real random variables \(X_i\), \(i\in I\), are independent
      exactly when, for every finite \(J\subseteq I\) and all bounded
      Borel functions \(f_i:\mathbb R\to\mathbb R\), \(i\in J\),
      \(\mathbb E\bigl[\prod_{i\in J}f_i(X_i)\bigr]
      =\prod_{i\in J}\mathbb E[f_i(X_i)]\).
\end{enumerate}
\end{hbtheorem}

\begin{hbtheorem}{thm:p15-w02b-p15-006-arbitrary-blocks}{proved=annex:p15-w02b-p15-006-arbitrary-blocks}{Disjoint groups of any size}
Let \(\mathcal F_i\), \(i\in I\), be independent sub-sigma-algebras
with \(I\) arbitrary, and let \((I_j)_{j\in L}\) be pairwise disjoint
subsets of \(I\) indexed by an arbitrary set \(L\). No group and no
number of groups is required to be countable. Put
\[
 \mathcal G_j\defeq\sigma\left(\bigcup_{i\in I_j}\mathcal F_i\right),
 \qquad j\in L,
\]
with \(\mathcal G_j=\{\varnothing,\Omega\}\) when \(I_j\) is empty.
Then the \(\mathcal G_j\) are independent, and any random elements
\(Y_j\) such that each \(Y_j\) is \(\mathcal G_j\)-measurable are
independent.
\end{hbtheorem}

\begin{hbexample}{def:p15-w02b-p15-006-mixed-family}{An uncountable independent family on the unit interval}
On the uniform probability space \((\Omega_U,\mathcal F_U,\mathbb U)\)
of \cref{thm:p15-lebesgue-unit-interval}, let \(Y_n\), \(n\geq1\), be
the independent fair bits of \cref{thm:p15-dyadic-independent-coins}.
For \(t\in(0,1)\), let \(W_t=\mathbf1_{\{t\}}\), the indicator of the
singleton event \(\{t\}\), which has probability zero. Index the
combined family by the disjoint union
\(I=\mathbb N_+\cup(0,1)\), with \(W_n=Y_n\) for \(n\in\mathbb N_+\).

The family \((W_i)_{i\in I}\) is independent. First replace every
\(W_t\), \(t\in(0,1)\), by the constant zero, and fix a finite
\(J\subseteq I\) and Borel sets \(B_i\), \(i\in J\). For a constant
index, \(\{0\in B_t\}\) is \(\Omega\) or \(\varnothing\). If any such
event is empty, both sides of the factorization are zero. Otherwise those
factors equal one and drop out, leaving the factorization for the fair
bits indexed by \(J\cap\mathbb N_+\), which holds by
\cref{thm:p15-dyadic-independent-coins}. The modified family is
therefore independent. Each \(W_t\) equals zero outside the null event
\(\{t\}\), so item~3 of \cref{thm:p15-w02b-p15-006-arbitrary-transfer}
transfers independence to \((W_i)_{i\in I}\).

Each \(W_t\) equals zero almost surely, yet at every outcome \(u\) the
member \(W_u\) equals one. Almost sure agreement of each member with a
constant therefore does not give almost sure agreement of the whole
uncountable family, and item~3 of
\cref{thm:p15-w02b-p15-006-arbitrary-transfer} needs none. The
uncountable part of this family is degenerate: each \(W_t\) is almost
surely constant. \Cref{thm:p15-w02b-p15-006-uniform-countable} shows that
on this space every independent family has this feature.
\end{hbexample}

\subsection{Uncountable families on a space with a countable approximating collection}
\label{sec:p15-w02b-p15-006-uncountable-limit}

\begin{hblemma}{thm:p15-w02b-p15-006-uniform-approximation}{proved=annex:p15-w02b-p15-006-uniform-approximation}{A countable approximating collection for the uniform space}
For an integer \(m\geq1\), call \((k/m,(k+1)/m]\), \(0\leq k<m\),
the cells of mesh \(1/m\). Let \(\mathcal D\) be the collection of all
sets \((0,1)\cap\bigcup_{k\in K}(k/m,(k+1)/m]\) with \(m\geq1\) and
\(K\subseteq\{0,\ldots,m-1\}\). Then \(\mathcal D\) is countable,
\(\mathcal D\subseteq\mathcal F_U\), and for every \(A\in\mathcal F_U\)
and every \(\varepsilon>0\) some \(D\in\mathcal D\) satisfies
\(\mathbb U(A\triangle D)<\varepsilon\), where
\(A\triangle D\defeq(A\setminus D)\cup(D\setminus A)\).
\end{hblemma}

\begin{hbtheorem}{thm:p15-w02b-p15-006-separable-limit}{proved=annex:p15-w02b-p15-006-separable-limit}{Nondegenerate independent events are countable on a separable space}
Suppose a countable collection \(\mathcal D\subseteq\mathcal F\) has
the property that for every \(A\in\mathcal F\) and every
\(\varepsilon>0\) some \(D\in\mathcal D\) satisfies
\(\mathbb P(A\triangle D)<\varepsilon\).
\begin{enumerate}
\item If \(A,B\) are independent events with
      \(\delta\leq\mathbb P(A),\mathbb P(B)\leq1-\delta\) for some
      \(\delta\in(0,1/2]\), then \(\mathbb P(A\triangle B)\geq2\delta^2\).
\item Every independent family of events \((A_i)_{i\in I}\) with
      \(\delta\leq\mathbb P(A_i)\leq1-\delta\) for all \(i\) and one
      \(\delta\in(0,1/2]\) has a countable index set.
\item In every independent family of events \((A_i)_{i\in I}\), the
      indices with \(0<\mathbb P(A_i)<1\) form a countable set.
\item In every independent family of real random variables
      \((X_i)_{i\in I}\), the indices at which \(X_i\) is not almost
      surely equal to a constant form a countable set. A real random
      variable \(X\) is almost surely constant exactly when
      \(\mathbb P(X\leq q)\in\{0,1\}\) for every rational \(q\).
\end{enumerate}
\end{hbtheorem}

\begin{hbcorollary}{thm:p15-w02b-p15-006-uniform-countable}{proved=annex:p15-w02b-p15-006-uniform-countable}{Independent families on the uniform space}
On the uniform probability space \((\Omega_U,\mathcal F_U,\mathbb U)\),
every independent family of events has only countably many members of
probability strictly between zero and one, and every independent family
of real random variables has only countably many members that are not
almost surely constant. An independent family with uncountably many
nondegenerate members therefore requires a probability space with no
countable approximating collection.
\end{hbcorollary}

\section{Pairwise independence, mutual independence and the complement test}
\label{sec:p15-w02b-p15-006-complement}

For finitely many events there are two equivalent ways to state mutual
independence. The subset form of
\cref{def:p15-finite-independence,def:p15-general-independence} imposes
one equation for each subset of the events. Grinstead and Snell give,
as an equivalent alternative, one equation for each way of choosing
every event or its complement and intersecting all of the choices
\cite[sec.~4.1, Definition~4.2, p.~141]{GrinsteadSnellChance2006}. In
the source each chosen event is written \(A'_j\), equal to \(A_j\) or
to its complement \(\tilde A_j\); here the complement is written
\(A_j^c\). For a proof the source refers only to the case of three
events, posed as its Exercise~33
\cite[sec.~4.1, p.~155]{GrinsteadSnellChance2006}.
The theorem below proves it for every \(n\), and extends it to
arbitrary families.

\begin{hbtheorem}{thm:p15-w02b-p15-006-complement-form}{proved=annex:p15-w02b-p15-006-complement-form}{The complement-choice form of mutual independence}
\par\nobreak
Let \(n\geq1\) and let \(A_1,\ldots,A_n\) be events. The following
are equivalent.
\begin{enumerate}
\item \emph{Subset form.} For every \(J\subseteq\{1,\ldots,n\}\),
      \(\mathbb P(\bigcap_{j\in J}A_j)=\prod_{j\in J}\mathbb P(A_j)\).
\item \emph{Complement-choice form.} For every choice
      \(B_j\in\{A_j,A_j^c\}\), \(1\leq j\leq n\),
      \[
       \mathbb P\left(\bigcap_{j=1}^nB_j\right)
                         =\prod_{j=1}^n\mathbb P(B_j).
      \]
\item \emph{Extended choice form.} For every choice
      \(C_j\in\{A_j,A_j^c,\Omega\}\), \(1\leq j\leq n\),
      \(\mathbb P(\bigcap_{j=1}^nC_j)=\prod_{j=1}^n\mathbb P(C_j)\).
\end{enumerate}
The second form consists of \(2^n\) equations over full \(n\)-fold
intersections. Neither the single equation with every \(B_j=A_j\) nor
the pairwise equations can replace it
(\cref{def:p15-w02b-p15-006-complement-examples}).
For an arbitrary family \((A_i)_{i\in I}\), the events are independent
exactly when the complement-choice form holds for every finite
subfamily, that is, for every finite \(J\subseteq I\) and every choice
\(B_i\in\{A_i,A_i^c\}\), \(i\in J\),
\(\mathbb P(\bigcap_{i\in J}B_i)=\prod_{i\in J}\mathbb P(B_i)\).
\end{hbtheorem}

\begin{hbexample}{def:p15-w02b-p15-006-complement-examples}{The complement test on two dependent triples}
\par\nobreak
\emph{A triple satisfying the full-intersection equation.} Grinstead
and Snell's Exercise~8 \cite[sec.~4.1, p.~151]{GrinsteadSnellChance2006}
takes \(\Omega=\{a,b,c,d,e,f\}\) with masses
\(m(a)=m(b)=1/8\) and \(m(c)=m(d)=m(e)=m(f)=3/16\), and the events
\(A=\{d,e,a\}\), \(B=\{c,e,a\}\), \(C=\{c,d,a\}\). Each event has
probability \(3/16+3/16+1/8=1/2\), and
\(\mathbb P(A\cap B\cap C)=m(a)=1/8=\mathbb P(A)\mathbb P(B)\mathbb P(C)\).
Every complement choice has product \(1/8\), but the eight full
intersections have probabilities
\[
\begin{array}{lclc}
 A\cap B\cap C=\{a\}:&1/8,&A\cap B^c\cap C^c=\varnothing:&0,\
 A\cap B\cap C^c=\{e\}:&3/16,&A^c\cap B\cap C^c=\varnothing:&0,\
 A\cap B^c\cap C=\{d\}:&3/16,&A^c\cap B^c\cap C=\varnothing:&0,\
 A^c\cap B\cap C=\{c\}:&3/16,&A^c\cap B^c\cap C^c=\{b,f\}:&5/16.
\end{array}
\]
They sum to one. Seven of the eight equations fail, so the complement
test rejects the triple although its one full-intersection equation
holds. The subset form also fails: \(\mathbb P(A\cap B)=m(e)+m(a)=5/16\)
differs from \(1/4\), and likewise for the other two pairs, so no two
of the events are independent.

\emph{A pairwise independent triple.} In the model of item~1 of
\cref{thm:p15-independence-counterexamples} (Durrett's Example~2.1.4
\cite[p.~45]{DurrettPTEAuthor2019}), with independent fair bits
\(X_1,X_2,X_3\) and \(A_1=\{X_2=X_3\}\), \(A_2=\{X_3=X_1\}\),
\(A_3=\{X_1=X_2\}\), every pair is independent. Any two of the
equalities imply the third, so a choice with exactly one complement
gives the empty intersection, of probability \(0\). With exactly two
complements, say \(A_1\cap A_2^c\cap A_3^c\), the bits satisfy
\(X_2=X_3\neq X_1\), which has two of the eight equally likely
outcomes and probability \(1/4\). Three complements require three
pairwise different bits, which is impossible. The eight probabilities
are therefore \(1/4\), three times \(0\), three times \(1/4\), and
\(0\), against the common product \(1/8\). The complement test rejects
the triple at every choice, although the pairwise equations all hold.
\end{hbexample}

\begin{hbwarning}{def:p15-w02b-p15-006-infinite-choice-warning}{Infinite intersections do not define independence}
\par\nobreak
For an infinite family, neither the full intersection nor the
complement-choice intersection over the whole index set can replace the
finite subfamilies. On the uniform space, with the fair bits \(Y_n\)
of \cref{thm:p15-dyadic-independent-coins}, put
\(A_1=A_2=\{Y_1=1\}\) and \(A_k=\{Y_k=1\}\) for \(k\geq3\). For any
choice \(B_k\in\{A_k,A_k^c\}\), \(k\geq1\), and any \(n\geq3\),
\[
 \mathbb U\left(\bigcap_{k\geq1}B_k\right)
   \leq\mathbb U\left(\bigcap_{k=3}^nB_k\right)=2^{-(n-2)},
\]
by monotonicity and the finite factorization of the fair bits. Hence
every such infinite intersection has probability zero, which equals
\(\lim_n\prod_{k=1}^n\mathbb U(B_k)=\lim_n2^{-n}\). Every infinite
choice equation of this kind therefore holds. Nevertheless
\(\mathbb U(A_1\cap A_2)=1/2\neq1/4=\mathbb U(A_1)\mathbb U(A_2)\),
so the family is not independent. Item~3 of
\cref{thm:p15-w02b-p15-006-finite-character} runs only from independence
to the infinite products, not back.
\end{hbwarning}

\begin{hbmethod}{def:p15-w02b-p15-006-method}{Testing independence of a family of any size}
Name the probability measure and the index set. For an infinite or
uncountable family, test each finite subfamily; an equation for the
intersection over the whole index set, with or without complements,
does not suffice. For finitely many events, either verify every subset
equation or verify all \(2^n\) complement-choice equations; checking
only the single \(n\)-fold intersection, or only pairs, does not
suffice. For sigma-algebras or random variables, verify the finite
factorizations on intersection-closed generating collections and apply
\cref{thm:p15-w02b-p15-006-arbitrary-transfer}; for grouped
information, apply \cref{thm:p15-w02b-p15-006-arbitrary-blocks}.
On the uniform space, or on any space with a countable approximating
collection, an asserted independent family with uncountably many
nondegenerate members is impossible by
\cref{thm:p15-w02b-p15-006-separable-limit}. The finite and countable
procedure of \cref{def:p15-independence-method} remains in force for
each finite subfamily.
\end{hbmethod}

\emph{Sources.} The definition through finite subcollections, the
finite definitions and the complement argument are from Durrett's
author version \cite[pp.~43--45, Theorems~2.1.1--2.1.3]{DurrettPTEAuthor2019};
the pi-system result applied to finite subfamilies is Durrett's
Theorem~2.1.7 \cite[p.~46]{DurrettPTEAuthor2019}. The complement-choice
form is Grinstead and Snell's Definition~4.2 with their Exercises~8
and~33 \cite[sec.~4.1, pp.~141, 151, 155]{GrinsteadSnellChance2006}.
The extension of the transfer and grouping results to uncountable
families, the infinite-choice boundary, and the countability theorem
for spaces with a countable approximating collection are proved here
and are not presented in either source.
